% Fixed135: General odd-index (rather than solely index three) Fox compiler
% for the marked Roe--Turturean G_Q2 source, a new index-five numerical
% source, strict cross-Sylow chain transitivity, explicit first Nielsen
% overlap homotopy, and a base-field coinduction corollary.

\subsection{Odd-index marked Schreier--Fox compilation and
strict cross-cover coherence}
\label{subsec:dyadic-f135-odd-sylow-compiler}

The fixed133--fixed134 strict Fox--Shapiro bridge
was proved by explicit Nielsen transformations for
one index-three nonnormal cover.  The structural
reason it works is more general: an odd-index
open subgroup of $G_{\mathbf Q_2}$ necessarily
contains the entire normal pro-$2$ wild
subgroup.  This makes the original
Roe--Turturean \emph{marked} quotient
universality stable under finite
wreath induction from such a subgroup.
The resulting pro-$2$ Schreier presentation
has exactly one more relation than its
augmentation-Fox rank.  Burnside--Nielsen
elimination then supplies a
group-level, integral finite-jet compiler,
without guessing a standard Demu\v{s}kin
word in advance.

Write
\[
\Gamma=G_{\mathbf Q_2},\qquad
W=\overline{\langle\!\langle
x_0,x_1\rangle\!\rangle}_{\Gamma},
\]
for the distinguished wild normal
pro-$2$ subgroup in the verified marked
Roe--Turturean source
\cite{RoeTurtureanGQ2}.
The four marked generators are
$(\sigma,\tau,x_0,x_1)$ and the
two relators are
$(r_t,r_w)$, where
\[
r_t=\sigma^{-1}\tau\sigma\tau^{-2},
\qquad
r_w=h_0u_1^{-1}x_1^\sigma c_0.
\]
As throughout the previous versions,
the requirement that $W$ be pro-$2$
is an \emph{additional marked
finite-quotient condition}, not
a consequence of the two naked words.
All statements below retain it.

\begin{lemma}[Odd index automatically absorbs wild inertia]
\label{lem:dyadic-f135-odd-index-wild-core}
Let $H\subset\Gamma$ be any open
subgroup of odd index
\[
m=[\Gamma:H].
\]
Then
\begin{equation}\label{eq:dyadic-f135-wild-in-core}
\boxed{\qquad
W\subseteq\operatorname{Core}_{\Gamma}(H)
=\bigcap_{g\in\Gamma}gHg^{-1}.
\qquad}
\end{equation}
In particular, in the finite
permutation action on the $m$
cosets $H\backslash\Gamma$,
the images of $x_0,x_1$
are trivial.  No normality of
$H$ is assumed.
\end{lemma}

\begin{proof}
Let $P$ be the finite image of
$\Gamma$ in the permutation group
of the cosets, and let $S\subset P$
be the stabilizer of the chosen
coset.  Its index is $m$, which
is odd.  Hence $S$ contains
a Sylow $2$-subgroup of $P$.
The image $\bar W$ of $W$ is
a normal finite $2$-subgroup
of $P$; every normal $2$-subgroup
is contained in every Sylow
$2$-subgroup.  Thus
$\bar W\subset S$.
Normality of $\bar W$ also gives
$\bar W\subseteq pSp^{-1}$
for every $p\in P$.
Therefore it lies in the core
of $S$, which is the kernel of
the permutation action.
Pulling back to $\Gamma$
proves \eqref{eq:dyadic-f135-wild-in-core}.
\end{proof}

\begin{theorem}[Uniform marked odd-index pro-$2$ Schreier presentation]
\label{thm:dyadic-f135-uniform-pro2-schreier}
Let $H\subset\Gamma$ be
open of odd index $m$,
and let $F/\mathbf Q_2$
be its finite fixed field,
so $[F:\mathbf Q_2]=m$.
Choose an ordered right
Schreier transversal
\[
\mathscr T=(t_0=1,t_1,\ldots,t_{m-1})
\quad\text{for }H\backslash\Gamma,
\]
written in the marked generators.
The indexed Schreier cover
of the \emph{two marked words}
$(r_t,r_w)$ has
\begin{equation}\label{eq:dyadic-f135-cover-cell-counts}
\boxed{
\begin{array}{c|ccc}
\text{cell degree}&0&1&2\\ \hline
\text{number}&m&4m&2m.
\end{array}}
\end{equation}
Contract a spanning tree
of its coset graph.
There remain
\begin{equation}\label{eq:dyadic-f135-cover-pro2-counts}
\boxed{\quad
d=3m+1\ \text{Schreier generators},
\qquad r=2m\ \text{indexed relators}.
\quad}
\end{equation}
Let $\mathscr P_H$ be the
pro-$2$ group defined by
these completed
Reidemeister--Schreier
relators, with
all profinite powers
rewritten in its
pro-$2$ quotients.
Then there is a
canonical isomorphism
of pro-$2$ groups
\begin{equation}\label{eq:dyadic-f135-marked-wreath-universality}
\boxed{\qquad
\mathscr P_H\simeq H(2)=G_F(2).
\qquad}
\end{equation}
It holds for
\emph{every} odd-index
$H$, tame or wild,
normal or nonnormal.
The assertion depends
on the original
\emph{marked wild}
pro-$2$ condition
of Roe--Turturean;
it does not treat
their two naked
full-group relators
as a complete
unmarked presentation.

The indexed
Schreier relator
words and their
completed Fox
derivatives admit a
terminating algorithm
at every fixed finite
pro-$2$ quotient or
coefficient-jet
precision.  The
algorithm is:
compute the finite
coset transition
table; perform
literal word rewriting
on its $m$ cosets;
replace each
profinite $\omega_2$
power by its
compatible finite
CRT exponent
in the chosen
finite wreath quotient;
then evaluate the
resulting words
in the finite
pro-$2$ stage.
\end{theorem}

\begin{proof}
Schreier's index
formula for an open
subgroup of the free
profinite group
on four generators
gives $1+m(4-1)=3m+1$
free generators after
contracting a coset
spanning tree.
The normal subgroup
generated by two
relators in the
ambient free profinite
group restricts to
the closed normal
subgroup generated
by their $m$ coset
conjugates.
This gives the cell
and relation counts
as formal Schreier
data.

The point is to
prove that the
additional \emph{marked}
wild condition induces
no unnoticed extra
relations in the
maximal pro-$2$
quotient of $H$.
Put $\mathcal F$ for the
free profinite group on
the four marked
letters, and let
$\mathcal E$ be the
preimage in $\mathcal F$
of $H$, an open
subgroup of index $m$.
The permutation
homomorphism from
$\mathcal F$ to its
finite coset-action
image $\Omega$ has
point stabilizer
corresponding to
$\mathcal E$.
By
Lemma~\ref{lem:dyadic-f135-odd-index-wild-core},
both wild letters
have trivial
permutation action.

Let $Q$ be an
arbitrary finite
$2$-group.
Every homomorphism
from $H$ to $Q$
annihilates the
two original
indexed relators;
therefore it gives
a homomorphism from
the formal pro-$2$
Schreier relator
quotient $\mathscr P_H$.
Conversely, start
with \emph{any}
homomorphism
$\psi:\mathcal E\to Q$
annihilating the
two sets of $m$
indexed relators.
The finite
wreath-induction
construction
associated to
the transversal
gives a continuous
homomorphism
\begin{equation}\label{eq:dyadic-f135-wreath-extension}
\mathcal F\longrightarrow
Q^m\rtimes\Omega,
\end{equation}
whose marked
relation values
vanish in every
coset coordinate:
this is exactly
what the indexed
Schreier relations
assert.
Both wild
generator images
belong to the
normal base group
$Q^m$, since
their permutation
parts are trivial.
Their full normal
closure in the
finite image is
therefore a
$2$-group.
Hence the induced
finite quotient is
\emph{admissible}
for the actual
Roe--Turturean
marked presentation.
It factors through
$\Gamma$, and its
restriction to $H$,
evaluated in the
distinguished coset
coordinate, is
the original $\psi$.

Thus the finite
$2$-group quotient
functors of
$\mathscr P_H$
and $H$ agree,
including their
marked source
maps.  Taking
inverse limits
proves
\eqref{eq:dyadic-f135-marked-wreath-universality}.
Finite CRT
representatives
of $\omega_2$
are compatible
in successive
finite wreath
quotients, giving
the stated
finite-stage
source-word and
Fox algorithm.
\end{proof}

\begin{theorem}[The universal $2m-1$ integral Nielsen pivot compiler]
\label{thm:dyadic-f135-primitive-pivot-compiler}
Retain an odd-index
$H=G_F\subset\Gamma$
as above.
The local maximal
pro-$2$ group
$P_F=H(2)$ is
a Demu\v{s}kin group
of minimal generator
rank
\begin{equation}\label{eq:dyadic-f135-local-rank}
\boxed{\qquad
d_{\min}=[F:\mathbf Q_2]+2=m+2,
\qquad r_{\min}=1.
\qquad}
\end{equation}
Let $J_{\mathrm{aug}}$
be the $2m\times(3m+1)$
\emph{mod-$2$
augmentation Fox matrix}
of the completed
Schreier relators from
Theorem~\ref{thm:dyadic-f135-uniform-pro2-schreier}.
Then
\begin{equation}\label{eq:dyadic-f135-fox-augmentation-rank}
\boxed{\qquad
\operatorname{rank}_{\mathbf F_2}
J_{\mathrm{aug}}
=(3m+1)-(m+2)=2m-1.
\qquad}
\end{equation}
In particular,
there is a
$(2m-1)\times(2m-1)$
Fox minor which is
a unit modulo $2$,
and it can be
\emph{found
deterministically}
by finite
$\mathbf F_2$
Gaussian elimination.

Choose $2m-1$
independent relators
and $2m-1$
original free
generators giving
such a unit minor.
Replace those
generators by
the selected
\emph{actual
completed relator
words}.
Together with
the remaining
$m+2$ generators
they form a
free pro-$2$ basis.
After this single
primitive Nielsen
change, delete the
$2m-1$ generator
letters which are
now defining
relations.  The
result is a genuine
one-relator minimal
presentation
\begin{equation}\label{eq:dyadic-f135-general-minimal-relation}
\boxed{\quad
P_F\cong
\langle z_1,\ldots,z_{m+2}
\mid\mathscr R_F\rangle_{\mathrm{pro}-2}.
\quad}
\end{equation}
The relator
$\mathscr R_F$ and
the inverse Nielsen
word map are
\emph{computable
at every finite
pro-$2$ quotient
and Zassenhaus
word jet}.
They need not be
single finite
ordinary words
or canonical under
a change of
Schreier transversal.

For every complete
finite projective
coefficient module
$B$ whose
$H$ action
factors through
$P_F$, these
genuine group
transformations
induce an
\emph{integral
Fox cochain}
equivalence
\begin{equation}\label{eq:dyadic-f135-uniform-cochain-cancellation}
\boxed{
\begin{aligned}
C_{\mathrm{cover}}^\bullet(H,B)
\cong {}&
C_{\mathrm{Fox}}^\bullet(P_F,B)\\
&\oplus
[B^{m-1}\xrightarrow{I}B^{m-1}]_{[0,1]}\\
&\oplus
[B^{2m-1}\xrightarrow{I}B^{2m-1}]_{[1,2]} .
\end{aligned}}
\end{equation}
The term
$C_{\mathrm{Fox}}^\bullet(P_F,B)$
is the genuine
$(1,m+2,1)$
completed Demu\v{s}kin
resolution.
The degree-wise
ranks of the
two sides of
\eqref{eq:dyadic-f135-uniform-cochain-cancellation}
are exactly
$(m,4m,2m)$.

Consequently the
indexed covering
cochain complex
computes
\begin{equation}\label{eq:dyadic-f135-uniform-schreier-rgamma}
\boxed{\quad
C_{\mathrm{cover}}^\bullet(H,B)
\simeq
R\Gamma(H,B)
\quad\text{in degrees }0,1,2,
\quad}
\end{equation}
compatibly with
complete derived
coefficient base
change and
all chosen
Fox presentation
transitions.
No global canonical
minimal basis is
required.
\end{theorem}

\begin{proof}
Since $F/\mathbf Q_2$
is finite and
$\mu_2\subset F$,
local Demu\v{s}kin
theory
\cite{NSW,LabuteClassification}
gives
$d(H(2))=[F:\mathbf Q_2]+2=m+2$
and one minimal
defining relation.
The mod-two trivial
coefficient differential
$d^0$ in any
pro-$2$ presentation
is zero; its
$H^1$ is the
annihilator of the
augmented relator
row matrix.
Theorem~\ref{thm:dyadic-f135-uniform-pro2-schreier}
identifies the
Schreier relator
quotient with
$H(2)$, so its
first-cohomology
dimension is
$m+2$.
Rank-nullity
proves
\eqref{eq:dyadic-f135-fox-augmentation-rank}.

Gaussian elimination
selects a unit
$(2m-1)$-minor.
The chosen
$2m-1$ relator
words have
linearly independent
images in the
Frattini quotient
of the free
pro-$2$ group
of rank $3m+1$.
The remaining
$m+2$ original
generators can
be selected to
complete those
vectors to a
basis.
The free
pro-$2$ Burnside
basis theorem
and Hopficity
then make the
corresponding
continuous
endomorphism
of the free
pro-$2$ group
an automorphism.
In its new
generators the
selected relators
are \emph{literally
generator letters}.
Killing those
letters with
their defining
relations leaves
the unique
unselected
relator, giving
\eqref{eq:dyadic-f135-general-minimal-relation}.
This is a
pro-$2$
Tietze--Nielsen
calculation, not
a mere Smith
equivalence of
numeric matrices.

For the Fox
cochain complex,
the two edges
of each selected
tree link provide
primitive
vertex--edge
unit differentials;
there are
$m-1$ of them.
The completed
Fox chain rule
sends each
pro-$2$ Nielsen
automorphism
to an invertible
completed
group-ring
Jacobian.
After the
selected relation
words have become
generator letters,
each relation
face has a
\emph{literal}
unit Fox derivative
in that coordinate.
Elementary
source, relation
and face
row operations
therefore split
the $2m-1$
degree-one--two
contractible
pairs, before
or after
evaluating
in $B$.
The remaining
three-term Fox
complex is
the minimal
one-relator
Demu\v{s}kin
complex.
Its exactness
over
$\mathbf Z_2[[P_F]]$
is the standard
pro-$2$
Poincar\'e
duality resolution.

Finally every
continuous pro-$2$
word computation
in the chosen
group can be
evaluated modulo
each finite
Zassenhaus
quotient.
The selected
basis automorphism
acts bijectively
on the compatible
finite quotient
system, so the
inverse can be
computed by
finite lifting
at each requested
precision.
The Fox
derivatives are
continuous and
the pivot matrices
invert over the
completed local
group algebra
because their
augmentation
determinants
are odd units.
This proves the
claimed
recursive effectivity
and derived
coefficient
compatibility.
\end{proof}

\begin{theorem}[Strict odd-index Fox--Shapiro completeness on every residual sector]
\label{thm:dyadic-f135-full-marked-all-residual-completeness}
Let $A$ be a complete
Noetherian local
$\mathbf Z_2$-algebra
with finite residue
field of characteristic
two, and let
$B$ be a finite
projective $A$-module
with continuous
$\Gamma=G_{\mathbf Q_2}$
action.
Let $\bar\rho_B(\Gamma)$
be its finite
residual image,
choose a Sylow
$2$-subgroup
$S$ of that
image, and set
\begin{equation}\label{eq:dyadic-f135-residual-sylow-preimage}
H=\bar\rho_B^{-1}(S),
\qquad
m=[\Gamma:H]
=[\bar\rho_B(\Gamma):S].
\end{equation}
Then $m$ is odd,
$W\subseteq H$,
and the
continuous $H$
action on $B$
factors through
the pro-$2$
group $H(2)$.

Let
\[
\mathcal J_\Gamma^\bullet(B):
\qquad
B\longrightarrow B^{4}
\longrightarrow B^2
\]
be the \emph{literal
Roe--Turturean
two-marked-relator}
coefficient Fox
complex,
with its
profinite $\omega_2$
powers evaluated
in the complete
coefficient
representation.
Put
\[
V=\operatorname{Coind}_H^\Gamma(B|_H)
\simeq B\otimes_{\mathbf Z_2}
\mathbf Z_2[H\backslash\Gamma].
\]
With a chosen
coset trivialization,
let
\begin{equation}\label{eq:dyadic-f135-odd-coset-projector}
\boxed{\quad
P_H=\frac1m\mathbf1_m\mathbf1_m^t,
\qquad
P_H^2=P_H.
\quad}
\end{equation}
Its induced
endomorphisms on
$V,V^4,V^2$
commute \emph{strictly}
with the two
Fox differentials,
and the
coinduced
Schreier complex
splits into
strict integral
cochain summands:
\begin{equation}\label{eq:dyadic-f135-uniform-coinduced-splitting}
\boxed{
\begin{aligned}
\mathcal J_\Gamma^\bullet(V)
&\cong
C_{\mathrm{cover}}^\bullet(H,B)\\
&\cong
\mathcal J_\Gamma^\bullet(B)
\oplus
\mathcal J_\Gamma^\bullet
\left(B\otimes\ker(\mathbf1_m^t)\right).
\end{aligned}}
\end{equation}
The first summand
is the image of
$P_H$ in every
cochain degree,
and represents the
actual derived
residual Sylow
projector
$\frac1m\operatorname{res}
\operatorname{cor}$.
For this \emph{arbitrary}
residual coefficient
representation
$B$ one therefore has
\begin{equation}\label{eq:dyadic-f135-universal-coefficient-rgamma}
\boxed{\qquad
\mathcal J_\Gamma^\bullet(B)
\simeq
R\Gamma(\Gamma,B)
\qquad
\text{in }D(A).
\qquad}
\end{equation}
The equivalence is
compatible with
complete derived
coefficient base
change and the
strict transitive
coset-projection
changes below.

This gives a
single fixed
\emph{four-generator,
two-marked-relator}
\emph{coefficient
compiler}
for \emph{every}
finite-projective
continuous
$2$-adic
representation
of $G_{\mathbf Q_2}$,
not only the
$S_3$ residual
sector of fixed134.
It does \emph{not}
assert that the
two naked relations
alone present
$G_{\mathbf Q_2}$,
or that their
formal two-face
group-algebra
complex is a
universal
finite free
$\mathbf Z_2[[\Gamma]]$
resolution independent
of the marked
wild condition.
\end{theorem}

\begin{proof}
The congruence
kernel of
$\mathrm{GL}_A(B)$
over its finite
residue image
is pro-$2$:
its successive
finite quotients
are additive
matrix groups
over residue
characteristic
two.
The residual
$H$ image is
the chosen
Sylow $2$-group.
Thus the full
$H$ image on
$B$ is an
extension of
pro-$2$ groups,
hence pro-$2$,
and factors through
$H(2)$.
The subgroup has
odd index
and contains
the normal wild
group by
Lemma~\ref{lem:dyadic-f135-odd-index-wild-core}.

Coinduction
is finite
because $H$
has finite
index.  A
given $\Gamma$
action on $B$
trivializes the
coinduced module
as the tensor
product of
$B$ and the
right-coset
permutation
module.
The constant
coset line and
augmentation
kernel split
over $A$ because
$m$ is a unit.
The constant
embedding and
the normalized
average are
$\Gamma$-equivariant,
so their idempotent
$P_H$ commutes
with all
completed Fox
coefficient
operations.
This proves
the strict
cochain splitting
in
\eqref{eq:dyadic-f135-uniform-coinduced-splitting}.

The left
coinduced
complex is
the actual
indexed
Schreier covering
cochain complex
by the
Fox--Shapiro
chain rule.
Theorems~\ref{thm:dyadic-f135-uniform-pro2-schreier}
and~\ref{thm:dyadic-f135-primitive-pivot-compiler}
show that it
is a finite
Fox resolution
of $H(2)$,
plus explicit
contractible
pairs.
Hence it
computes
$R\Gamma(H,B)$.
Shapiro's lemma
identifies the
latter with
$R\Gamma(\Gamma,V)$.
The constant
direct summand
of the coinduced
object corresponds
to the ordinary
restriction map
$R\Gamma(\Gamma,B)
\to R\Gamma(H,B)$,
while the
sum map corresponds
to corestriction.
Consequently its
normalized
projector is
$\frac1m
\operatorname{res}
\operatorname{cor}$.
This is an
integral
idempotent with
image
$R\Gamma(\Gamma,B)$
since
$\operatorname{cor}
\operatorname{res}=m$.

The same
idempotent
splits the
\emph{literal
marked coefficient
complex}
$\mathcal J_\Gamma^\bullet(V)$
as
$\mathcal J_\Gamma^\bullet(B)$
plus the augmentation
factor.
The image
is therefore
identified
with
$R\Gamma(\Gamma,B)$,
proving
\eqref{eq:dyadic-f135-universal-coefficient-rgamma}.
Every step is
defined over
the complete
coefficient ring
using only odd
index inverses,
pro-$2$ Fox
Nielsen transformations
and completed
coefficient
base change.
\end{proof}

\begin{corollary}[One fixed marked $\mathbf Q_2$ compiler for every dyadic local field]
\label{cor:dyadic-f135-all-finite-extensions-coinduced}
Let $K/\mathbf Q_2$
be \emph{any} finite
extension, not
necessarily of odd
degree or Galois.
Let $A$ and
$B$ be as in
Theorem~\ref{thm:dyadic-f135-full-marked-all-residual-completeness},
but let $B$ carry a
continuous $G_K$
action.
Put
\[
n=[K:\mathbf Q_2],
\qquad
V_K=\operatorname{Coind}_{G_K}^{G_{\mathbf Q_2}}B.
\]
Then $V_K$
is finite projective
over $A$ of rank
$n\operatorname{rank}_A B$,
and the same fixed
Roe--Turturean
four-generator,
two-marked-relator
cochain formulas
provide
\begin{equation}\label{eq:dyadic-f135-all-k-via-q2-marked}
\boxed{\quad
\mathcal J_{G_{\mathbf Q_2}}^\bullet(V_K)
\simeq
R\Gamma(G_{\mathbf Q_2},V_K)
\simeq
R\Gamma(G_K,B).
\quad}
\end{equation}
Thus the fixed
marked source
computes \emph{every}
dyadic local
coefficient cohomology
object in three
finite cochain
degrees, without
asserting a
uniform explicit
four-generator
word presentation
of each
$G_K$.
The local-field
coset data and
the induced
transport matrices
are extra finite
inputs, not a
new intrinsic
group presentation
of $G_K$.
\end{corollary}

\begin{proof}
Finite-index
coinduction from
$G_K$ to
$G_{\mathbf Q_2}$
preserves finite
projective
coefficient modules
and continuity.
Apply
Theorem~\ref{thm:dyadic-f135-full-marked-all-residual-completeness}
to $V_K$,
then the
continuous
finite-index
Shapiro lemma
to recover
$R\Gamma(G_K,B)$.
The index
$n$ need
not be odd:
the odd-index
Sylow restriction
is selected
\emph{after}
the coefficient
module $V_K$
has been
coinduced to
$G_{\mathbf Q_2}$.
\end{proof}

\begin{theorem}[Strict independence of different odd-index Fox covers]
\label{thm:dyadic-f135-cross-sylow-strict-coherence}
Fix a complete
residual coefficient
module $B$ as above,
and let
$H_1,H_2,\ldots$
be \emph{any}
odd-index open
subgroups on which
$B$ has pro-$2$
image.
For each $H_i$,
write
\[
C_i^\bullet
=C_{\mathrm{cover}}^\bullet(H_i,B),
\quad
\iota_i:\mathcal J_\Gamma^\bullet(B)\to C_i^\bullet,
\quad
p_i:C_i^\bullet\to\mathcal J_\Gamma^\bullet(B)
\]
for the strict
constant-coset
embedding and
normalized average,
so $p_i\iota_i=I$.
Define the
marked-to-marked
comparison for
two different
covers by
\begin{equation}\label{eq:dyadic-f135-cover-change-operations}
\boxed{\qquad
\Phi_{ji}:=\iota_jp_i:
C_i^\bullet\longrightarrow C_j^\bullet.
\qquad}
\end{equation}
Then \emph{on the
nose as cochain maps}
\begin{equation}\label{eq:dyadic-f135-strict-transitive-odd-covers}
\boxed{
\begin{aligned}
\Phi_{kj}\Phi_{ji}&=\Phi_{ki},\\
\Phi_{ij}\Phi_{ji}
&=\iota_i p_i=:E_i,\\
E_i^2&=E_i.
\end{aligned}}
\end{equation}
Thus the projected
direct summands
$\operatorname{im}E_i$
are canonically
strictly isomorphic,
and \emph{all}
higher overlap
associators for
these projected
coefficient
carriers are
identity maps.
No intersection
$H_i\cap H_j$
is required;
indeed, such an
intersection
need not have
odd index.

There is also a
first explicit
\emph{minimal-Fox
overlap homotopy}.
Choose for each
$i$ the integral
Fox cancellations of
Theorem~\ref{thm:dyadic-f135-primitive-pivot-compiler},
as a strong
deformation
retract
\[
F_i^\bullet
\ \substack{\xrightarrow{\ j_i\ }\\[-1ex]
\xleftarrow[\ q_i\ ]{}}\ 
C_i^\bullet,
\qquad
q_i j_i=I,\quad
I-j_iq_i=d h_i+h_i d,
\]
where $F_i$ is the
chosen minimal
Demu\v{s}kin
Fox complex
for $H_i$,
including the
appropriate
coefficient module.
Define
\begin{equation}\label{eq:dyadic-f135-minimal-fox-overlap-maps}
T_{ji}:=q_j\Phi_{ji}j_i:
F_i^\bullet\longrightarrow F_j^\bullet.
\end{equation}
For every triple
the explicitly
given degree-minus-one
map
\begin{equation}\label{eq:dyadic-f135-first-fox-overlap-homotopy}
\boxed{\quad
\mathfrak H_{kji}
=q_k\Phi_{kj}h_j\Phi_{ji}j_i
\quad}
\end{equation}
satisfies
\begin{equation}\label{eq:dyadic-f135-minimal-fox-homotopy-identity}
\boxed{\quad
T_{ki}-T_{kj}T_{ji}
=d\mathfrak H_{kji}
+\mathfrak H_{kji}d.
\quad}
\end{equation}
All such
comparisons and
homotopies are
integral,
base-change
compatible on
the chosen
finite Fox gauges,
and recursively
computable at
every finite
coefficient/word
precision.
The strict
projector layer
requires \emph{no}
higher fillers;
at the chosen
minimal-Fox
gauges the
displayed $\mathfrak H$
is the first
nontrivial
coherence correction.
The complete
all-arity
noncanonical
minimal-Fox
higher-coherence
calculus remains
separate from
the strict
projector assertion.
\end{theorem}

\begin{proof}
All $\iota_i$
and $p_i$
are strict
cochain maps,
and $p_i\iota_i=I$.
Thus
\[
\Phi_{kj}\Phi_{ji}
=\iota_k(p_j\iota_j)p_i
=\iota_kp_i=\Phi_{ki}.
\]
The same
multiplication
gives the
idempotent and
inverse
identities.
These are
exact equalities
of matrices
on cochain
degrees $0,1,2$,
independent of
the odd indices
and transversals.
No common
subgroup is used.

For the minimal
Fox gauges,
the actual
Tietze/Schreier
unit cancellations
give strong
deformation
retracts with
the displayed
homotopies.
Their overlap
maps compose as
\[
T_{kj}T_{ji}
=q_k\Phi_{kj}
(j_jq_j)\Phi_{ji}j_i.
\]
Substitute
$j_jq_j=I-dh_j-h_jd$.
Strict
transitivity of
$\Phi$ and
the cochain-map
identities of
$q,j,\Phi$
then show
\[
T_{ki}-T_{kj}T_{ji}
=d(q_k\Phi_{kj}h_j\Phi_{ji}j_i)
+(q_k\Phi_{kj}h_j\Phi_{ji}j_i)d,
\]
which is
\eqref{eq:dyadic-f135-minimal-fox-homotopy-identity}.
Each $h_i$
is constructed
using inverse
unit Fox pivots
and free
pro-$2$ basis
changes;
its finite-stage
matrix entries
are therefore
effective and
$2$-integral.
\end{proof}

\begin{theorem}[A new independent index-five tame source certificate]
\label{thm:dyadic-f135-index-five-exact-certificate}
Put
\[
F_5=\mathbf Q_2(\sqrt[5]2),
\quad
L_5=F_5(\zeta_5),
\quad
\Gamma=G_{\mathbf Q_2}.
\]
The extension $F_5/\mathbf Q_2$
is totally tamely ramified
of degree five, while
$\zeta_5$ generates the
unramified degree-four
extension.  Its Galois
closure has the
non-pro-$2$ tame
quotient
\begin{equation}\label{eq:dyadic-f135-index-five-quotient}
\boxed{\quad
\operatorname{Gal}(L_5/\mathbf Q_2)
\simeq C_5\rtimes_{2}C_4,
\qquad
[\,\Gamma:G_{F_5}\,]=5.
\quad}
\end{equation}
The associated
Sylow preimage
$H_5=G_{F_5}$
has a three-term
Schreier cover
with ranks
$(5,20,10)$,
and its minimal
pro-$2$
Demu\v{s}kin quotient
has seven generators
and one relator.

Let $\Sigma,\Upsilon\in
M_5(\mathbf Z_2)$
be the exact
permutation matrices
for the coset
actions
\[
\Sigma_{ij}
=\delta_{j,\,2i\bmod5},
\qquad
\Upsilon_{ij}
=\delta_{j,\,i+1\bmod5},
\quad
0\le i,j<5.
\]
Then
$\Sigma^{-1}\Upsilon\Sigma=\Upsilon^2$,
$\Sigma^4=\Upsilon^5=I_5$.
Put
\begin{equation}\label{eq:dyadic-f135-index-five-norm}
N_5=\sum_{j=0}^4\Upsilon^j,
\qquad
P_5=\frac15N_5
=\frac15\mathbf1_5\mathbf1_5^t .
\end{equation}
On the coinduced
trivial coefficient
module $\mathbf Z_2^5$,
the \emph{literal}
marked full-source
degree-one Fox row is
\begin{equation}\label{eq:dyadic-f135-index-five-fox-blocks}
\boxed{
\begin{aligned}
d^1_{\mathrm{tame}}(a,b)
&=\Sigma^{-1}(\Upsilon-I)a+
   (\Sigma^{-1}-I-\Upsilon)b,\\
d^1_{\mathrm{wild}}(b,c,d)
&=(\Sigma^{-2}+2I)P_5b\\
&\quad+
\bigl((\Sigma^{-2}+3I)P_5-2I\bigr)c\\
&\quad+
(\Sigma^{-1}-P_5)d .
\end{aligned}}
\end{equation}
After multiplying
the five wild
face rows by
the odd $2$-adic
unit $5$,
the resulting
\emph{integer}
$10\times20$
matrix is
\begin{equation}\label{eq:dyadic-f135-index-five-scaled-jacobian}
\boxed{
\mathsf J_5=
\begin{pmatrix}
\Sigma^{-1}(\Upsilon-I)&
\Sigma^{-1}-I-\Upsilon&
0&0\\
0&(\Sigma^{-2}+2I)N_5&
(\Sigma^{-2}+3I)N_5-10I&
5\Sigma^{-1}-N_5
\end{pmatrix}.
}
\end{equation}
It has
$\mathbf F_2$-rank nine.
Two explicit
integer determinant
witnesses, with
rows and columns
numbered from one,
are
\begin{equation}\label{eq:dyadic-f135-index-five-minors}
\boxed{
\begin{aligned}
\det(\mathsf J_5[
\{1,\dots,9\},
\{1,2,3,4,6,16,17,18,19\}])
&=-125,\\
\det(\mathsf J_5[
\{1,\dots,10\},
\{1,2,3,4,6,16,17,18,19,11\}])
&=-1250 .
\end{aligned}}
\end{equation}
Thus the complete
$\mathbf Z_2$-Smith
invariant list is
\begin{equation}\label{eq:dyadic-f135-index-five-smith}
\boxed{\qquad
\operatorname{SNF}_{\mathbf Z_2}
(\mathsf J_5)
=(1,1,1,1,1,1,1,1,1,2).
\qquad}
\end{equation}
The ten face
rows possess
exactly nine
primitive
relation--edge
directions and one
$2$-primary top
class.  By the
general primitive
Nielsen compiler,
the actual
minimal group
presentation has
seven generators
and one relation,
with
\[
H^2(G_{F_5},\mathbf Z_2)
\cong\mathbf Z/2\mathbf Z.
\]
This is the first
independent higher-index
test of the general
odd-index theorem,
not a restatement
of the index-three
$S_3$ case.
\end{theorem}

\begin{proof}
The polynomial
$X^5-2$ is
Eisenstein at
two.  The
multiplicative
order of $2$
modulo $5$
is four, so
$\mathbf Q_2(\zeta_5)$
is unramified
of degree four,
disjoint from
the totally ramified
degree-five
extension.
This gives
the tame metacyclic
Galois group and
the index-five
Sylow complement.

In the affine
semidirect extension
of the five-sheet
permutation module
by the tame quotient,
the element $\Sigma$
has $2$-power
residual order four,
so its $\omega_2$
projection is itself.
The element
$\Upsilon$ has
odd residual order
five, and the
$\omega_2$
projection of
$(\Upsilon,v)$
is the pure
translation
$(I,P_5v)$.
In particular,
for the marked
wild letter $x_j$
with residual
permutation identity,
\[
(x_j\tau)^{\omega_2}
=(I,P_5(c_j+b))
\]
in affine coordinates.
The auxiliary wild
words then give
\[
d_0=P_5b+(P_5-I)c,\qquad
g_0=\sigma_2^2
\quad\text{with
linear part }\Sigma^2.
\]
The remaining
commutators of
pure translations
vanish.
The wild
derivative is
therefore
\[
(\Sigma^{-2}+I)c+
(\Sigma^{-2}+3I)d_0
-P_5(b+d)
+\Sigma^{-1}d,
\]
which simplifies
to
\eqref{eq:dyadic-f135-index-five-fox-blocks}.
The tame
derivative follows
from
$r_t=\sigma^{-1}\tau
\sigma\tau^{-2}$.
Multiplying
the wild rows by
five gives
\eqref{eq:dyadic-f135-index-five-scaled-jacobian}.

Finite
$\mathbf F_2$
row reduction
gives rank nine.
Computing the
two displayed
minors by
integer
Gaussian/Bareiss
elimination gives
$-125$ and
$-1250$.
The first is odd,
so the first nine
Smith factors
are units.
Every full
$10\times10$
minor is even
because the
mod-two rank
is nine;
the second
displayed minor
has valuation
one.
Hence the final
Smith factor is
exactly two,
as claimed.
The same conclusion
follows from
Theorem~\ref{thm:dyadic-f135-primitive-pivot-compiler}
and the local
$q(F_5)=2$
Demu\v{s}kin
classification,
but the two
integer witnesses
supply an
independent
finite algebra
audit.
\end{proof}

\begin{corollary}[Further tame indices and the sharp remaining boundary]
\label{cor:dyadic-f135-tame-indices-and-scope}
For any odd
integer $m>1$
such that
the multiplicative
order
$f=\operatorname{ord}_m(2)$
is a power of two,
the tame
Galois closure
of
$\mathbf Q_2(\sqrt[m]2)$
has quotient
\[
C_m\rtimes_2 C_f
\]
with a Sylow-$2$
complement of
index $m$.
The odd-index
Fox compiler
applies to its
degree-$m$
field with
cell ranks
$(m,4m,2m)$,
minimal Demu\v{s}kin
generator rank
$m+2$, and
$2m-1$
unit relation
directions.
The integer
index examples
$m=3,5,15,17$
have been
independently
checked through
modulo-four
unit-pivot
elimination:
their augmented
$2m\times4m$
Fox matrices
have mod-two
rank $2m-1$
and one surviving
Smith factor
of valuation one.

The group-level
compiler and
odd-index
constant-coset
projector are
therefore no
longer restricted
to the $S_3$
residual quotient.
What remains
\emph{strictly
stronger} is
a single
presentation-independent
minimal Fox
basis functor
with canonical
all-arity
higher homotopies
for general
dyadic fields.
The present
version gives
coefficientwise
uniform
cohomological
completeness
and strict
coherent
\emph{projected}
cover changes,
with an explicit
first
minimal-Fox
overlap homotopy;
it does not
assert a global
closed ordinary-word
minimal gauge.
\end{corollary}

\begin{proof}
The tame
quotient is the
semidirect
group generated
by an odd-order
inertia element
$\tau$ of order
$m$ and a
Frobenius element
$\sigma$ of
order
$f=\operatorname{ord}_m(2)$,
with
$\sigma^{-1}\tau
\sigma=\tau^2$.
When $f$ is a
power of two,
the Frobenius
complement is
a Sylow-$2$
subgroup and
has odd index $m$.
Apply
Theorems~\ref{thm:dyadic-f135-uniform-pro2-schreier},
\ref{thm:dyadic-f135-primitive-pivot-compiler},
and
\ref{thm:dyadic-f135-full-marked-all-residual-completeness}.
The four
finite matrix
checks are
recorded in the
reproducible
modulo-four
unit-pivot
certificate
for this
version.
The difference
between
strict projected
coherence and
canonical all-arity
minimal gauge
is explained by
Theorem~\ref{thm:dyadic-f135-cross-sylow-strict-coherence}.
\end{proof}

\begin{warning}[The marked wild condition remains essential]
\label{warn:dyadic-f135-not-naked-full-relators}
The two full-group
Roe--Turturean
words $(r_t,r_w)$
are \emph{not}
assumed to
present the
absolute Galois
group as a bare
unmarked
two-relator
profinite quotient.
The crucial
step in the
uniform odd-index
proof is the
finite wreath
argument, which
shows that
induced finite
$2$-group
quotients are
admissible
\emph{because}
the wild
generator images
lie in the normal
$2$-group base.
Without this
fact there is
no reason for
the $2m$
Schreier word
relations alone
to present $H(2)$.
Similarly, the
coefficientwise
completeness
of the literal
two-face Fox
complex is not
a claim of a
fixed full-group
projective
$\mathbf Z_2[[\Gamma]]$
resolution
independent of
residual-sector
and marked
normal-closure
input.

The equality
of projected
odd-index
cochains is
\emph{strict},
but the
minimal-Fox
Nielsen gauges
are chosen,
not canonical.
The displayed
triple-overlap
homotopy is
a concrete
finite-level
coherence step,
not a proof of
every higher
arity in one
preferred
closed form.
\end{warning}
